\section{Consistency Boundaries}
\label{sec:ConsistencyBoundaries}

This section explains what the Raft implementation guarantees when the
network partitions and the cluster appears to split into disconnected
groups. It separates two concerns:

\begin{enumerate}
\item consistency of the replicated Raft log and the state machine driven
by that log---Raft addresses this directly,
\item consistency of everything around that log---clients, retries, caches,
projections, and external side effects---which remains an application and
integration responsibility.
\end{enumerate}

\subsection{What ``split brain'' means here}

In informal distributed-systems language, ``split brain'' usually means a
network partition where the cluster is divided into separate sets of nodes
where more than one set of nodes behaves as if it is the authoritative cluster.
In a ``split brain'' scenario scenario, what we fear is that writes could proceed 
independently in individual sets, each believing they are \emph{the} cluster, and 
then even after healing the data sets would diverge. 

A correct Raft system, guarantees that this outcome can not happen for 
committed state.

What \emph{can} happen under a partition:

\begin{itemize}
\item The cluster partitions into a majority side and one or more minority
sides.
\item Nodes on multiple sides may temporarily believe they should lead.
\item An old leader may continue acting as leader until it notices quorum
loss.
\item Clients may see redirects, retries, timeouts, or ambiguous outcomes.
\end{itemize}

What should \emph{not} happen:

\begin{itemize}
\item Two different partitions both committing conflicting log entries for
the same log position in the same configuration.
\item A minority partition continuing to commit authoritative writes as if
it still had quorum.
\end{itemize}

The useful distinction is: apparent split brain at the transport or
liveness level can happen; committed split brain at the replicated-log
level should not.

\subsection{Why ``majority'' and ``minority'' are well-defined}

The terms majority and minority are not decided dynamically by ``who can
currently see whom.'' In Raft, elections and commitment are evaluated
against the configured cluster membership---a specific current voting
set. A candidate wins only by receiving votes from a majority of that
configured set; an entry commits only when replicated to a majority of
that configured set.

If a 5-node voting configuration \texttt{\{A, B, C, D, E\}} has quorum
threshold 3, and the network partitions into \texttt{\{A, B, C\}} and
\texttt{\{D, E\}}, then \texttt{\{A, B, C\}} is the majority side and
\texttt{\{D, E\}} is the minority side. This holds regardless of whether
\texttt{D} and \texttt{E} can still reach each other. Reachability affects
liveness; configuration defines legitimacy.

The same logic applies to any partition shape. A 7-node cluster split
into 4 + 3: only the 4 side can make committed progress because quorum
is 4. The cluster cannot legitimately shrink itself by accident just
because nodes disappeared behind a partition. A new voting set only becomes
authoritative after the cluster commits that configuration change through
the log.

\subsection{What Raft guarantees}

For the authoritative replicated state machine, Raft gives these safety
properties:

\begin{itemize}
\item Only a leader with a quorum can commit new log entries.
\item Once an entry is committed, future leaders must contain that entry.
\item Conflicting uncommitted entries on losing sides are eventually
overwritten.
\item After partition healing, followers converge to the majority-selected
log.
\item Membership changes are serialized through the log and protected by
joint consensus.
\end{itemize}

Applied to this implementation:

\begin{itemize}
\item Ordinary writes are accepted only through the leader path.
\item Write success is defined as committed and applied, not merely appended
locally.
\item Linearizable reads are served only by a leader that can establish a
fresh quorum-backed read window.
\item Removed nodes decommission themselves after finalized removal.
\end{itemize}

In practice, under a partition: the majority side may continue to elect a
leader and commit writes; the minority side may be alive but should lose
the ability to commit new authoritative writes; once the partition heals,
minority-side uncommitted divergence is discarded.

The Jepsen partition scenarios (Section~\ref{sec:JepsenTesting}) validate
this: isolating one node, isolating the leader, cutting the leader into a
minority, and overlapping reconfiguration with a partition---all checking
that only one side keeps making committed progress.

\subsection{What Raft does not guarantee}

Raft does not automatically make the rest of the application ecosystem
consistent. It does not solve:

\begin{itemize}
\item client retries after ambiguous timeouts,
\item duplicate command submission,
\item side effects that happen before commit is known,
\item cache invalidation across failover,
\item stale read models or projections,
\item reconciliation of data written outside the Raft leader path,
\item domain-level merge policy for concurrent offline edits,
\item idempotency of calls into external systems.
\end{itemize}

This is the point often mislabeled as ``eventual consistency.'' Inside the
Raft state machine, the goal is strong consistency for committed state.
Around the state machine, many components are only eventually consistent:
materialized views, denormalized read stores, caches, search indexes,
downstream analytics pipelines, webhook consumers, external billing or
provisioning systems. Those must converge eventually, but Raft alone does
not define how.

\subsection{The three outcome classes}

When a client issues a command during instability, there are three
fundamentally different outcomes:

\begin{description}
\item[Definitely committed] The leader had quorum, the log entry committed,
and the state machine applied it. The clean case.
\item[Definitely not committed] The request was rejected because the node
was not leader, the cluster had no usable leader, the request was invalid
or unauthorized, or the CAS precondition did not match. Safe to treat as no
state change.
\item[Ambiguous] The client timed out, disconnected, or lost contact while
the cluster was unstable. The command may have committed, or it may have
been dropped. The client cannot safely assume either outcome.
\end{description}

The ambiguity exists even in a correct Raft system because the client can
lose the response after the cluster has already committed the command. The
Jepsen harness treats such outcomes as indeterminate rather than definite
failures for this reason.

\subsection{Where inconsistency can still appear}

\subsubsection{Uncommitted local divergence}

A minority or superseded leader may append entries locally that never
commit. That local log tail is inconsistent with the surviving majority
for a while. Raft handles this by rollback and overwrite during
reconciliation. Application state changes driven only by committed entries
do not need custom business logic for this case.

\subsubsection{Side effects emitted too early}

If application code sends an email, charges a card, provisions a resource,
or updates another datastore before the Raft command is known to be
committed, a partition or failover may produce an external side effect
with no corresponding committed entry in the authoritative log. Raft
cannot repair that automatically.

\subsubsection{Duplicate execution from retries}

If a client times out and retries without an idempotency key, the cluster
may eventually apply the same logical intent twice. Raft preserves log
consistency, but not business-level uniqueness unless the application
models it.

\subsubsection{Stale or lagging read models}

If a system derives secondary views from the Raft log asynchronously,
those views may lag after partitions, failovers, or restart recovery.
That is not a Raft safety failure; it is an eventually consistent
projection.

\subsubsection{Writes outside the leader path}

If any service allows local writes on followers, offline replicas, edge
nodes, or separate databases not serialized through the Raft leader, then
conflict resolution becomes an application concern. At that point you are
no longer relying on Raft alone for authoritative writes.

\subsection{What application developers must accommodate}

\subsubsection{Treat authoritative state as log-derived}

The cleanest model: only committed Raft log entries mutate authoritative
state. Everything else is derived from or synchronized to that state.
This prevents most split-brain reconciliation problems.

\subsubsection{Use idempotency keys for client commands}

Clients should attach a stable command identity for operations that may
be retried. This lets the application distinguish ``same logical command
retried after timeout'' from ``new logical command.'' Without that,
ambiguous client failures can become duplicate business actions.

\subsubsection{Delay irreversible side effects until commit is known}

If a side effect is not reversible, it should not happen before the
corresponding Raft command is committed and applied. Safer patterns:
commit first, then emit side effects from an outbox or committed-event
stream; or emit only compensatable side effects.

\subsubsection{Make side-effect consumers replay-safe}

Anything downstream of the committed log should tolerate duplicate
delivery, delayed delivery, and reordered observation across independent
streams. This is standard outbox and projection discipline.

\subsubsection{Define reconciliation rules for external projections}

For caches, search indexes, or read models, developers need to define how
staleness is detected, how a projection catches up after outage, whether
rebuild-from-log is supported, and whether consumers should prefer
monotonic versions or log indexes.

\subsubsection{Define domain merge rules if offline mutation is allowed}

If the product intentionally allows updates outside the single Raft leader
path, then Raft no longer provides the whole answer. Developers must define
last-writer-wins or version-based conflict rules, per-field merge behavior,
tombstone/delete precedence, and human review workflows for irreconcilable
conflicts.

\subsection{Recommended documentation for applications}

Any real application using this Raft layer should document:

\begin{enumerate}
\item which data is authoritative and log-backed,
\item which APIs return definite success, definite failure, or ambiguous
outcome,
\item how client retries are deduplicated,
\item when side effects are emitted relative to commit,
\item how projections rebuild or catch up after outage,
\item whether any writes occur outside the Raft leader path,
\item what conflict resolution policy applies if they do.
\end{enumerate}

If those items are left implicit, teams often assume Raft solves more than
it actually does.

\subsection{Bottom line}

Raft prevents committed split brain for the authoritative replicated log.
Raft does not eliminate ambiguous client outcomes during failures,
duplicate intent from retries, stale projections, premature external side
effects, or domain conflicts created outside the leader/commit path.

The right mental model: Raft gives strong consistency for committed state;
everything around that state still needs explicit application-level
consistency design.
